% RESUMO--------------------------------------------------------------------------------

\begin{resumo}[RESUMO]
\begin{SingleSpacing}

% Não altere esta seção do texto--------------------------------------------------------
\imprimirautorcitacao. \imprimirtitulo. \imprimirdata. \pageref {LastPage} f. \imprimirprojeto\ – \imprimirprograma, \imprimirinstituicao. \imprimirlocal, \imprimirdata.\\
%---------------------------------------------------------------------------------------

Este relatório apresenta as atividades desenvolvidas no Laboratório Assert do Instituto Federal da Paraíba (IFPB), Campus João Pessoa, entre fevereiro e junho de 2026, por meio do PIAGI. O estágio teve como objetivo desenvolver um módulo para o sistema Banco de Especialistas, visando formalizar e rastrear solicitações de alterações cadastrais e de classificação dos usuários. Para isso, implementou-se o módulo de Solicitações, contemplando a modelagem das entidades Solicitação e Snapshot, as regras de negócio para criação, edição, aprovação e rejeição de solicitações e uma linha do tempo (timeline) para registrar os eventos e as decisões do processo. O desenvolvimento também abrangeu a implementação de operações de persistência e endpoints, a validação e serialização de dados, a integração com a arquitetura existente e a construção de funcionalidades de interface para solicitantes e avaliadores. A solução utilizou Python, Flask, PostgreSQL, SQLAlchemy e Marshmallow no back-end, Nuxt no front-end e Docker e Docker Compose para configuração do ambiente. Também foram desenvolvidos testes automatizados com Pytest para verificar as funcionalidades implementadas. Como resultado, o módulo integrou o fluxo de solicitações à plataforma, reduzindo a dependência de comunicações por e-mail e intervenções manuais no banco de dados, além de proporcionar maior formalização, rastreabilidade e histórico das alterações realizadas.\\

\textbf{Palavras-chave}: Desenvolvimento full-stack. Banco de Especialistas. Solicitações. Rastreabilidade. Testes automatizados.

\end{SingleSpacing}
\end{resumo}

% OBSERVAÇÕES---------------------------------------------------------------------------
% Altere o texto inserindo o Resumo do seu trabalho.
% Escolha de 3 a 5 palavras ou termos que descrevam bem o seu trabalho 

